
\uline{un nombre premier} @@M0@@ \uline{et} @@M1@@ \uline{est un groupe de torsion. Alors toute application polynôme} @@M2@@ \uline{qui se factorise par} @@M3@@\uline{, est constante.}

\emph{Démonstration} : Il s'agit de montrer que pour tout @@M4@@, on a @@M5@@.

a) Supposons d'abord que @@M6@@ soit un élément de torsion, disons @@M7@@ (@@M8@@, @@M9@@). Alors le polynôme @@M10@@ a un nombre infini de racines, à savoir tous les entiers divisibles par @@M11@@. Donc le polynôme @@M12@@ est identiquement nul et on a @@M13@@.

b) Supposons maintenant que @@M14@@ soit un élément libre de @@M15@@ divisible par toutes les puissances @@M16@@ (@@M17@@). Considérons des polynômes @@M18@@ définis par
@@M19@@
Pour tout @@M20@@ on a l'identité
@@M21@@
i.e.
@@M22@@
d'où on tire par comparaison de coefficients
@@M23@@
